% !TeX program = lualatex
% Compile twice: lualatex open_problems_500.tex
% All references and prose are embedded; no BibTeX database or external inputs.
\documentclass[11pt,a4paper]{article}
\usepackage[margin=24mm,headheight=14pt]{geometry}
\usepackage{fontspec}
\setmainfont{Latin Modern Roman}
\setsansfont{Latin Modern Sans}
\setmonofont{Latin Modern Mono}
\usepackage{amsmath,amssymb,mathtools}
\usepackage{unicode-math}
\setmathfont{Latin Modern Math}
\usepackage{microtype}
\usepackage{enumitem,xurl,needspace}
\usepackage[dvipsnames]{xcolor}
\usepackage{hyperref}
\hypersetup{unicode=true,colorlinks=true,linkcolor=MidnightBlue,urlcolor=MidnightBlue,
 pdftitle={500 Mathematical Problem Statements},pdfauthor={Source-annotated compilation},bookmarksnumbered=true}
\usepackage{bookmark}
\usepackage{tocloft}
\setlength{\cftsecnumwidth}{3em}
\cftsetpnumwidth{2.5em}
\cftsetrmarg{4em}
\usepackage{titlesec}
\titleformat{\section}{\Large\bfseries\raggedright}{\thesection}{0.7em}{}
\usepackage{fancyhdr}
\pagestyle{fancy}\fancyhf{}
\fancyhead[L]{\small\sffamily Mathematical problem statements}
\fancyhead[R]{\small Source edition 22}
\fancyfoot[C]{\thepage}
\setlength{\parindent}{0pt}
\setlength{\parskip}{5pt plus 1pt}
\setlength{\emergencystretch}{3em}
\setcounter{tocdepth}{1}
\setcounter{secnumdepth}{1}
\setlist[itemize]{leftmargin=3em,itemsep=5pt,topsep=4pt,parsep=0pt}
\allowdisplaybreaks
\newcommand{\N}{\mathbb{N}}
\newcommand{\Z}{\mathbb{Z}}
\newcommand{\Q}{\mathbb{Q}}
\newcommand{\R}{\mathbb{R}}
\newcommand{\C}{\mathbb{C}}
\newcommand{\F}{\mathbb{F}}
\newcommand{\A}{\mathbb{A}}
\newcommand{\PP}{\mathbb P}
\newcommand{\E}{\mathbb{E}}
\newcommand{\eps}{\varepsilon}
\newcommand{\dd}{\,\mathrm d}
\newcommand{\sref}[2]{\hyperlink{source:#1:#2}{[S#2]}}
\newcommand{\eref}[2]{\hyperlink{extra:#1:#2}{[E#2]}}
% Turn the next switch off for a shorter reading copy. The quotations preserve
% every exactTarget field, including qualifications omitted from a short summary.
\newif\ifshowcatalogue
\showcataloguetrue
\newcommand{\cataloguescope}[1]{\ifshowcatalogue
 \paragraph{Exact scope in the supplied catalogue.}
 \begin{quote}\small #1\end{quote}\fi}

% Standalone research notebook, original catalogue problem 188.
% Compile twice with: lualatex 188.tex
\numberwithin{equation}{section}
\hypersetup{pdftitle={Research attempt -- problem 188},pdfauthor={Research notebook; original compilation retained}}
\fancyhead[L]{\small\sffamily Research notebook}
\fancyhead[R]{\small\sffamily Problem 188 | 27 September 2026}
\begin{document}
\setcounter{section}{187}
% ================================================================
% problemId: problem.collatz-conjecture
% source releaseRank: 188; source releaseStatus: open
% exactTarget SHA-256: 5608318e4ea73247cfb12c609de785f3d73c444e86e855d622aee9051aa5c51e
\clearpage
\section{\texorpdfstring{Collatz conjecture}{Collatz conjecture}}\label{188}
\subsection{Definitions and mathematical statement}
Define $T:\N_{>0}\to\N_{>0}$ by
\[
 T(n)=\begin{cases}n/2,&2\mid n,\\3n+1,&2\nmid n.\end{cases}
\]
The Collatz assertion is $\forall n\ge1\ \exists k\ge0$ such that $T^k(n)=1$, where $T^0$ is the identity. Reaching one enters the cycle $1,4,2,1$. The iteration here uses the unaccelerated map; no division by an additional power of two is included in an odd step. A negative resolution may be an orbit entering a different cycle or one that never repeats and never reaches one. Verifying any finite initial range does not decide the universal assertion.
\par\smallskip\textit{Formulation and concept references:} \sref{188}{1}.
\cataloguescope{Prove or disprove that, for every positive integer N, repeated application of N -> N/2 when N is even and N -> 3N+1 when N is odd eventually reaches 1.}
\subsection{Short English statement}
Does repeatedly halving even integers and replacing odd integers by three times the number plus one always eventually reach one?
% BEGIN RESEARCH ADDITION ATTEMPT-188
\subsection{Research attempt}

\paragraph{Current frontier.}
Tao's theorem says that almost all starting values, in logarithmic density, have an orbit reaching a bound tending arbitrarily slowly to infinity. It does not say that every orbit reaches one; the quantifier distinction survives in the revised manuscript \sref{188}{1}. An exact finite computation cannot remove the exceptional set. We pursue a deterministic descent certificate and test whether adding residue information to a logarithmic potential can make it monotone. The resulting obstruction is proved below and does not depend on statistical modeling of parity bits.

\paragraph{Attack route.}
One possible proof would establish that every odd $n>1$ eventually descends below $n$, after which strong induction finishes the conjecture. Another would construct a well-founded potential. A natural computable candidate is $V(n)=c\log n+f(n)$, where $f$ is a bounded correction, possibly depending on finitely many residues. A third route is a finite cover by variable-length descent certificates; its missing requirement is coverage of all exceptional residue cylinders, not just almost every cylinder. We test the potential route first, including fixed-block variants.

\paragraph{Attempt.}
For odd positive $n$ define
\[
 U(n)=\frac{3n+1}{2^{v_2(3n+1)}}.
\]
The original unaccelerated orbit reaches one if and only if its odd-to-odd orbit under $U$ does. The omitted intermediate integers are exactly the forced successive halvings, so no additional orbit is introduced; $U(1)=1$ represents the original $1,4,2,1$ cycle.

\textbf{Lemma (arbitrarily long expanding words).} For integers $L>R\ge1$, set $n_0=2^L-1$. Then
\begin{equation}\label{eq188:rise}
 U^j(n_0)=3^j2^{L-j}-1\quad(0\le j\le R),
\end{equation}
and all the first $R$ valuations equal one.
Indeed, if $n_j=3^j2^{L-j}-1$, then
$3n_j+1=2(3^{j+1}2^{L-j-1}-1)$. For $j<R<L$ the parenthesis is odd. This proves the formula by induction. Every step in this segment increases the odd value.

\textbf{Proposition (bounded-correction obstruction).} There do not exist $c>0$ and a bounded function $f$ on the positive odd integers for which
\begin{equation}\label{eq188:impossible}
 V(U(n))\le V(n)\quad\text{for every odd }n>1,
 \qquad V(n)=c\log n+f(n).
\end{equation}
Suppose $|f|\le M$. Telescoping \eqref{eq188:impossible} along \eqref{eq188:rise} gives
\[
 c\log\frac{3^R2^{L-R}-1}{2^L-1}\le2M.
\]
First choose $R$ with $cR\log(3/2)>2M$, and then let $L\to\infty$. The left side tends to $cR\log(3/2)$, a contradiction. This rules out even nonincrease, hence also strict decrease. It applies to every bounded correction, not only periodic corrections.

A fixed block length does not repair the obstruction. If
$V(U^r(n))\le V(n)$ holds for all sufficiently large odd $n$, apply it at $n_0,n_r,\ldots,n_{(j-1)r}$ along an expanding word of length $jr$. Letting $j$ grow again contradicts boundedness of $f$. For a residue correction there is an even more direct test: choose
\[
 n_0=2^{r+1}mt-1\quad(t\ge1).
\]
The first $r$ iterates all satisfy $n_j\equiv-1\pmod m$, so a function depending only on $n\bmod m$ cancels exactly between the endpoints while the logarithm strictly increases.

We next record the exact affine formula behind a more flexible residue search. Prescribe positive integers $a_1,\ldots,a_r$, and put $A_j=a_1+\cdots+a_j$, $A_0=0$. Whenever the first $r$ valuations equal that word,
\begin{equation}\label{eq188:word}
 U^r(n)=\frac{3^rn+B_r}{2^{A_r}},\qquad
 B_r=\sum_{j=0}^{r-1}3^{r-1-j}2^{A_j}.
\end{equation}
This follows from $2^{a_{j+1}}n_{j+1}=3n_j+1$. The numerator congruence requiring the final quotient to be odd singles out one odd residue modulo $2^{A_r+1}$; reversing the recurrence shows that it gives exactly the prescribed valuations. Thus the word has density $2^{-A_r}$ among odd integers. This is a finite-cylinder statement, not independence along one infinite integer orbit.

For such a word a certified descent occurs whenever
\[
 2^{A_r}>3^r,\qquad
 n>\frac{B_r}{2^{A_r}-3^r}.
\]
The affine offset matters: a negative logarithmic drift alone does not imply descent for every integer in the cylinder. A proposed proof by residue certificates must cover every odd integer greater than a verified finite base range with some word meeting this inequality. We have not obtained such a cover. Formula~\eqref{eq188:rise} proves that no fixed maximum word length can give a descent certificate for all large inputs.

\paragraph{Stress test.}
The expanding words are finite. They neither produce a divergent orbit nor a nontrivial cycle: after the prescribed word, a much larger power of two may divide the next iterate. Our obstruction therefore does not refute Collatz. It refutes a specific universal monotonicity requirement, even though that requirement initially looks compatible with average negative drift. Conversely, almost-everywhere stopping results do not rule out an exceptional orbit. A minimal counterexample would have to avoid every verified descent cylinder, but no compactness principle says a residual set of $2$-adic cylinders contains, or does not contain, a positive integer with the required infinite trajectory.

The checking script verifies long exact expanding segments, the affine word formula, and the single-residue assertion for bounded word lengths. None of these computations is described as checking the entire conjecture. We make no claim that the obstruction lemma has not appeared elsewhere.

\paragraph{Outcome.}
\textbf{Outcome: Exploratory approach with unresolved gap.}
The proposed bounded-correction logarithmic Lyapunov strategy fails, and the failure has a precise arbitrary-length construction. Exact variable-word descent certificates remain viable, but a proof must allow unbounded word lengths or a genuinely unbounded correction and must establish universal coverage. That missing deterministic control, rather than a finite-range computation or an average drift calculation, is what would be required to resolve the target.

% END RESEARCH ADDITION ATTEMPT-188
\subsection{Sources}
\begin{itemize}\raggedright
\item[\textnormal{[S1]}]\hypertarget{source:188:1}{} Terence Tao, Forum of Mathematics, Pi 10 (2022), e12.
\url{https://arxiv.org/abs/1909.03562}.
{\small\textit{Catalogue formulation source.}}
\end{itemize}

\end{document}
